\chapter{Rank and the Fundamental Theorem of Linear Algebra}

A single integer captures the true geometric dimensionality of a matrix
transformation: the \textbf{rank}. A matrix may contain thousands of rows and
columns yet collapse the entire input space onto a single line, or it may be
small and fill the entire space it maps into. The dimensions $m \times n$
describe the size of the matrix array, but the rank measures the dimension of
its image. This chapter explores what rank means, why the row rank always equals
the column rank, and how rank organizes the four fundamental subspaces.

We also show that every matrix decomposes into a sum of rank-one outer
products, providing the conceptual foundation for the singular value
decomposition.

\chapterline

\section{Rank is the true size of a matrix}

Let us start with a concrete geometric example. Take the matrix
\[
  \mat{A} = \begin{bmatrix} 1 & 2 \\ 2 & 4 \end{bmatrix}.
\]
It is a $2\times 2$ matrix. But watch what it does to vectors. Feed it $\colvec{1 \\ 0}$ and you
get $\colvec{1 \\ 2}$. Feed it $\colvec{0 \\ 1}$ and you get $\colvec{2 \\ 4}$.
Both outputs sit on the exact same line through the origin -- the line in the direction
$\colvec{1 \\ 2}$, because $\colvec{2 \\ 4} = 2\colvec{1 \\ 2}$. In fact
\emph{every} input comes out on that one line:
\[
  \mat{A}\colvec{x \\ y} = \colvec{x + 2y \\ 2x + 4y} = (x + 2y)\colvec{1 \\ 2}.
\]
The entire input plane $\R^2$ maps onto a single line. Although represented by
a $2 \times 2$ array of numbers, the transformation is fundamentally
one-dimensional.

\begin{intuitionbox}
The \textbf{rank} of a matrix is the dimension of the space its outputs fill --
the dimension of the column space $\col(\mat{A})$. In plain terms, it is the
number of independent directions the matrix can actually reach. A $1000 \times 1000$
matrix whose output always lands on a single line has rank $1$. The dimensions
$m \times n$ count the entries; the rank counts the dimension of the column
space.
\end{intuitionbox}

\begin{figure}[H]
\centering
\begin{tikzpicture}[scale=1.0,>=Stealth]
  % left: input plane with grid
  \begin{scope}[shift={(-3.6,0)}]
    \draw[help lines,gray!40] (-1.6,-1.6) grid (1.6,1.6);
    \draw[->,gray] (-1.8,0)--(1.8,0);
    \draw[->,gray] (0,-1.8)--(0,1.8);
    % nullspace line crushed to 0
    \draw[targetgreen,thick,dashed] (-1.6,0.8)--(1.6,-0.8);
    \node[targetgreen,font=\scriptsize] at (1.1,-1.15) {$\nullsp(\mat{A})\colon x+2y=0$};
    \draw[->,vecblue,very thick] (0,0)--(1,0) node[below right]{$\vec{e}_1$};
    \draw[->,vecred,very thick] (0,0)--(0,1) node[above left]{$\vec{e}_2$};
    \node at (0,-2.2) {\small input plane $\R^2$};
  \end{scope}
  % arrow
  \draw[->,thick] (-1.5,0)--(-0.3,0) node[midway,above]{$\mat{A}$};
  % right: collapsed onto a line
  \begin{scope}[shift={(2.4,0)}]
    \draw[->,gray] (-1.8,0)--(1.8,0);
    \draw[->,gray] (0,-1.8)--(0,2.1);
    \draw[structurepurple,very thick,dashed] (-0.9,-1.8)--(1.1,2.2);
    \draw[->,vecblue,very thick] (0,0)--(0.5,1.0) node[below right=-1pt, font=\footnotesize]{$\mat{A}\vec{e}_1$};
    \draw[->,vecred,very thick] (0,0)--(1.0,2.0) node[right=2pt, font=\footnotesize]{$\mat{A}\vec{e}_2 = 2\mat{A}\vec{e}_1$};
    \node[structurepurple, anchor=east, font=\small] at (-0.95,-1.55) {$\col(\mat{A})$};
    \node at (0,-2.25) {\small output: one line};
  \end{scope}
\end{tikzpicture}
\caption{The matrix $\protect\mat{A}=\protect\begin{bmatrix}1&2\\2&4\protect\end{bmatrix}$ takes the full
two-dimensional input plane and crushes it onto a single line. Its rank is $1$, while the dashed line
$x+2y=0$ is crushed completely to the origin.}
\label{ch16:collapse}
\end{figure}

Compare this with
\[
  \mat{B} = \begin{bmatrix} 1 & 2 \\ 2 & 5 \end{bmatrix}.
\]
One number changed -- the bottom-right went from $4$ to $5$ -- and the columns
are now linearly independent. Now $\mat{B}\colvec{1\\0} = \colvec{1\\2}$ and
$\mat{B}\colvec{0\\1} = \colvec{2\\5}$, and those two outputs do \emph{not} lie
on a common line: $\colvec{2\\5}$ is not a multiple of $\colvec{1\\2}$. The two
columns point in different directions, so together they span all of $\R^2$. The
rank is $2$. A single entry change alters the rank from $1$ to $2$, transforming
a singular matrix into an invertible transformation that spans all of $\R^2$.

\begin{focusbox}
Three equivalent ways to say what rank is. Hold all three; they are the same
fact seen from three sides.
\begin{itemize}
  \item \textbf{Column rank:} the number of linearly independent columns, i.e.
        $\dimn \col(\mat{A})$.
  \item \textbf{Row rank:} the number of linearly independent rows, i.e.
        $\dimn \row(\mat{A})$.
  \item \textbf{Pivot count:} the number of pivots you find when you row-reduce
        the matrix to echelon form.
\end{itemize}
The deep theorem of this chapter is that these three numbers are \emph{always
equal}. We call the common value $\rank(\mat{A})$.
\end{focusbox}

\subsection{Reading rank off the echelon form}

You already know elimination. It is the most reliable rank detector we have.
Reduce a matrix to row echelon form and count the pivots -- the leading nonzero
entries, one per pivot row. That count is the rank.

\begin{example}\label{ch16:pivot-count}
Find the rank of
\[
  \mat{C} = \begin{bmatrix} 1 & 2 & 3 \\ 2 & 4 & 8 \\ 3 & 6 & 7 \end{bmatrix}.
\]
Eliminate below the first pivot. Subtract $2\times$row$_1$ from row$_2$ and
$3\times$row$_1$ from row$_3$:
\[
  \begin{bmatrix} 1 & 2 & 3 \\ 0 & 0 & 2 \\ 0 & 0 & -2 \end{bmatrix}.
\]
The second column never produced a pivot -- the $(2,2)$ position is zero and so
is everything below it. Move to column $3$: there is a pivot in row $2$. Add
row$_2$ to row$_3$ to clear the last entry:
\[
  \begin{bmatrix} 1 & 2 & 3 \\ 0 & 0 & 2 \\ 0 & 0 & 0 \end{bmatrix}.
\]
Two pivots (in columns $1$ and $3$). So $\rank(\mat{C}) = 2$. The third row was
a combination of the first two all along (specifically,
$\text{row}_3 = 5\,\text{row}_1 - \text{row}_2$); elimination exposed it as a
row of zeros.
\end{example}

\begin{trapbox}
Rank is \emph{not} ``the number of nonzero rows you started with,'' and it is not
``the number of columns minus the number of zeros.'' A matrix can have no zero
entries anywhere and still have rank $1$ (every row a multiple of the first).
Rank is about \emph{independence}, not about where the zeros are. You must do the
elimination -- or otherwise check for dependence -- to know it.
\end{trapbox}

\begin{labbox}
Compute the rank of each by hand, reducing to echelon form and counting pivots:
\[
  \begin{bmatrix} 2 & 4 \\ 1 & 2 \end{bmatrix}, \qquad
  \begin{bmatrix} 1 & 0 & 1 \\ 0 & 1 & 1 \\ 1 & 1 & 2 \end{bmatrix}, \qquad
  \begin{bmatrix} 1 & 2 & 1 \\ 1 & 3 & 2 \\ 1 & 4 & 3 \end{bmatrix}.
\]
(Answers: $1$, $2$, $2$. In the second, row$_3 = $ row$_1+$row$_2$; in the third,
row$_3 = 2\,$row$_2 - $row$_1$. Hunt for those dependencies before you reduce --
seeing them is the central skill.)
\end{labbox}

\section{Rank-nullity: the conservation law}

Here is the single most useful counting fact in linear algebra. When a matrix
acts on $\R^n$, every one of the $n$ input dimensions has to go \emph{somewhere}.
Some of them survive and become independent output directions -- those are
counted by the rank. The rest get crushed to zero -- those are counted by the
dimension of the nullspace, $\dimn\nullsp(\mat{A})$, the number we call the
\textbf{nullity}. (In plain terms, nullity is the dimension of the matrix's blind
spot: the number of independent directions that collapse completely to the zero
vector.) Nothing else can happen to a dimension. So the two counts must add up
to the number of input dimensions you started with.

\begin{theorem}[Rank--Nullity]\label{ch16:rank-nullity}
For any $m \times n$ matrix $\mat{A}$ (with real entries),
\[
  \rank(\mat{A}) + \dimn\nullsp(\mat{A}) = n,
\]
where $n$ is the number of columns -- the dimension of the input space.
\end{theorem}

\begin{intuitionbox}
Rank--nullity expresses an exact dimensional balance for any linear mapping
$\mat{A}: \R^n \to \R^m$:
\begin{center}
$\underbrace{n}_{\text{input dimension}} \;=\;
 \underbrace{\rank(\mat{A})}_{\text{dimension of image}} \;+\;
 \underbrace{\dimn\nullsp(\mat{A})}_{\text{dimension of kernel}}.$
\end{center}
The total dimension $n$ of the domain is partitioned between the directions that
contribute to the output (the row space) and the directions mapped to zero (the
nullspace). Increasing the rank necessarily decreases the nullity by the same
amount.
\end{intuitionbox}

\subsection{Why it is true, from the echelon form}

The cleanest proof rides on something you already trust: every column of $\mat{A}$
either holds a pivot or it does not. Reduce $\mat{A}$ to reduced row echelon form
$\mat{R}$. The columns split into two kinds.

\begin{itemize}
  \item \textbf{Pivot columns:} there are $r = \rank(\mat{A})$ of them. The
        corresponding columns of the \emph{original} $\mat{A}$ are linearly
        independent and form a basis for the column space.
  \item \textbf{Free columns:} there are $n - r$ of them, one for each free
        variable. Each free variable produces exactly one special solution of
        $\mat{A}\vec{x} = \vec{0}$, and those special solutions form a basis for
        the nullspace.
\end{itemize}

So the nullspace has dimension $n - r$. That is the entire theorem:
$r + (n - r) = n$. Let us make the special-solution count airtight, because that
is the part that does the work.

\begin{proof}
Row reduction does not change the nullspace: $\mat{A}\vec{x}=\vec0$ and
$\mat{R}\vec{x}=\vec0$ have exactly the same solutions, because each elimination
step is a reversible operation on the equations. So it suffices to count the
nullspace dimension of $\mat{R}$.

Say $\mat{R}$ has $r$ pivots, hence $r$ pivot columns and $n-r$ free columns.
Each free column corresponds to a free variable that we are allowed to set to
anything. Form the $j$-th special solution by setting the $j$-th free variable
to $1$, all other free variables to $0$, and solving for the pivot variables from
the equations $\mat{R}\vec{x}=\vec0$ (each pivot variable is determined by the
free ones). This gives $n-r$ vectors $\vec{s}_1,\dots,\vec{s}_{n-r}$.

\emph{They are independent.} Look only at the free-variable positions. There the
vectors $\vec{s}_1,\dots,\vec{s}_{n-r}$ read off as the standard basis vectors of
$\R^{n-r}$ -- a single $1$ in distinct slots. A combination
$c_1\vec{s}_1+\cdots+c_{n-r}\vec{s}_{n-r}=\vec0$ forces every $c_j=0$ just by
reading those slots. So the special solutions are linearly independent.

\emph{They span the nullspace.} Take any $\vec{x}$ with $\mat{R}\vec{x}=\vec0$.
Let its free-variable values be $c_1,\dots,c_{n-r}$. The vector
$\vec{x} - (c_1\vec{s}_1+\cdots+c_{n-r}\vec{s}_{n-r})$ is also in the nullspace
(the nullspace is closed under combinations), and it has \emph{all} free
variables equal to zero. But once the free variables are zero, the equations
$\mat{R}\vec x = \vec 0$ force every pivot variable to zero too. So that
difference is $\vec0$, meaning $\vec{x}$ \emph{is} the combination. The special
solutions span.

Independent and spanning: the $n-r$ special solutions are a basis, so
$\dimn\nullsp(\mat{A}) = n-r$. Adding $r$ to both sides gives $r + (n-r) = n$.
\end{proof}

\begin{example}\label{ch16:rn-example}
Take a $2\times 3$ matrix, already in a convenient form:
\[
  \mat{A} = \begin{bmatrix} 1 & 0 & 2 \\ 0 & 1 & 3 \end{bmatrix}.
\]
Here $n = 3$ columns. Pivots in columns $1$ and $2$, so $r = \rank(\mat{A}) = 2$.
Column $3$ is free. Rank--nullity predicts $\dimn\nullsp = 3 - 2 = 1$: a single
line of solutions. Let us find it. Set the free variable $x_3 = 1$; the equations
$x_1 + 2x_3 = 0$ and $x_2 + 3x_3 = 0$ give $x_1 = -2$, $x_2 = -3$. The one special
solution is
\[
  \vec{s} = \colvec{-2 \\ -3 \\ 1}, \qquad
  \nullsp(\mat{A}) = \spn\{\vec{s}\}, \quad \dimn = 1. \checkmark
\]
Two input dimensions produce nonzero output (the rank), one was mapped to
zero (the nullity), and $2 + 1 = 3$. The dimension sum matches the theorem.
\end{example}

\begin{walkbox}
A $5\times 5$ matrix has rank $3$. Without computing anything: how many dimensions
does it map to zero? How many independent directions can its outputs reach?
Could a different $5\times 5$ matrix have rank $3$ and map only one dimension
to zero? Walk until you are certain the answer to the last question is ``no,
never'' -- and you can state clearly which dimension theorem guarantees this.
\end{walkbox}

\begin{strangbox}
Strang emphasizes this dimension balance throughout his lectures, writing
$r + (n-r) = n$ as the core identity connecting rank and nullity. In geometric
terms, the dimension of the domain equals the dimension of the image plus the
dimension of the kernel: $\dimn(\text{domain}) = \dimn(\text{image}) +
\dimn(\text{kernel})$. This is the rank--nullity theorem for linear maps, which
serves as the linear algebra formulation of the first isomorphism theorem.
\end{strangbox}

\section{Row rank equals column rank}

We have seen that ``number of independent columns'' and ``number of independent
rows'' both equal the number of pivots. Because the columns live in $\R^m$ and
the rows live in $\R^n$, it is not immediately obvious that the maximum number
of linearly independent columns must always equal the maximum number of
linearly independent rows. Yet this equality holds for every matrix.

\begin{theorem}[Row rank $=$ Column rank]\label{ch16:row-col}
For every matrix $\mat{A}$, the dimension of the column space equals the dimension
of the row space.
\end{theorem}

\subsection{The decisive picture: $\mat{A}$ as a short product}

Here is the argument I find most convincing, because it is constructive and you
can see every piece. Let $r$ be the column rank of $\mat{A}$ ($m\times n$). Pick a
basis $\vec{c}_1,\dots,\vec{c}_r$ for the column space -- $r$ vectors in $\R^m$.
Gather them as the columns of an $m \times r$ matrix $\mat{C}$.

Every column $\vec{a}_j$ of $\mat{A}$ lives in the column space, so it is some
combination of the basis: $\vec{a}_j = \mat{C}\vec{f}_j$ for a coefficient vector
$\vec{f}_j \in \R^r$. Stack those coefficient vectors as columns of an $r\times n$
matrix $\mat{F}$. Then, column by column,
\[
  \mat{A} = \mat{C}\,\mat{F}, \qquad
  \underbrace{\mat{A}}_{m\times n} =
  \underbrace{\mat{C}}_{m\times r}\,
  \underbrace{\mat{F}}_{r\times n}.
\]
This is the \textbf{rank factorization}: every matrix of column rank $r$ is a
``tall times wide'' product through a waist of width $r$. (Figure
\ref{ch16:rankfac}.) Now read off the rows. In the product $\mat{C}\mat{F}$, every
row of $\mat{A}$ is a combination of the $r$ rows of $\mat{F}$ -- because row $i$
of $\mat{A}$ equals (row $i$ of $\mat{C}$) $\times \mat{F}$, a weighted sum of
$\mat{F}$'s rows. So the row space of $\mat{A}$ is spanned by the $r$ rows of
$\mat{F}$, which means
\[
  \text{row rank of } \mat{A} \;\le\; r \;=\; \text{column rank of } \mat{A}.
\]
That is half the theorem. For the other half, apply this same inequality to the
transpose $\mat{A}\T$. The rows of $\mat{A}$ are the columns of $\mat{A}\T$ and
vice versa, so the inequality for $\mat{A}\T$ reads
$\text{column rank of }\mat{A} \le \text{row rank of }\mat{A}$. Two inequalities
pointing opposite ways force equality. Row rank $=$ column rank. $\qquad\blacksquare$

\begin{figure}[H]
\centering
\begin{tikzpicture}[scale=0.95]
  \def\mh{2.2} % matrix height
  % A
  \draw[fill=sourceblue!12,draw=sourceblue,thick] (0,0) rectangle (2.6,\mh);
  \node at (1.3,\mh/2) {$\mat{A}$};
  \node[below=2pt] at (1.3,0) {\small $m\times n$};
  \node at (3.3,\mh/2) {$=$};
  % C (tall, narrow)
  \draw[fill=targetgreen!14,draw=targetgreen,thick] (4.0,0) rectangle (4.9,\mh);
  \node at (4.45,\mh/2) {$\mat{C}$};
  \node[below=2pt] at (4.45,0) {\small $m\times r$};
  % F (short, wide)
  \draw[fill=quotientorange!16,draw=quotientorange,thick] (5.3,\mh/2-0.45) rectangle (7.9,\mh/2+0.45);
  \node at (6.6,\mh/2) {$\mat{F}$};
  \node[below=2pt] at (6.6,\mh/2-0.45) {\small $r\times n$};
  % waist annotations
  \draw[<->,structurepurple,thick] (4.0,-0.75)--(4.9,-0.75);
  \node[structurepurple,below=1pt] at (4.45,-0.75) {\small $r$ cols};
  \draw[<->,structurepurple,thick] (8.15,\mh/2-0.45)--(8.15,\mh/2+0.45);
  \node[structurepurple,right=1pt] at (8.15,\mh/2) {\small $r$ rows};
  \node[structurepurple,font=\small\bfseries] at (6.6,-0.75) {waist $= r$};
\end{tikzpicture}
\caption{The rank factorization $\protect\mat{A}=\protect\mat{C}\protect\mat{F}$.
Every column of $\protect\mat{A}$ is built from the $r$ columns of $\protect\mat{C}$
(so column rank $\le r$); every row of $\protect\mat{A}$ is built from the $r$ rows
of $\protect\mat{F}$ (so row rank $\le r$). The single number $r$ -- the width of
the waist -- governs both.}
\label{ch16:rankfac}
\end{figure}

\begin{aestheticbox}
Although columns live in $\R^m$ and rows live in $\R^n$, both counts reflect
the same structural rank. In the rank factorization $\mat{A} = \mat{C}\mat{F}$,
the intermediate dimension $r$ is the minimum inner dimension required to factor
$\mat{A}$. This factorization demonstrates that the row rank and column rank
are simply two perspectives on the same intrinsic invariant of the linear
transformation.
\end{aestheticbox}

\begin{labbox}
\textbf{Mental check: the bottleneck in action.}
If $\mat{A}$ is a $4 \times 100$ matrix, what is the largest possible column rank
it can have? What is the largest possible row rank?
(Answer: At most $4$ for both! Although there are a hundred columns, they live in
$\R^4$, so at most four can be linearly independent. And because row rank equals
column rank, those four dimensions create a waist of width at most $4$, preventing
the hundred rows from ever spanning more than a four-dimensional space.)
\end{labbox}

\begin{historybox}
In the nineteenth century, mathematicians studying systems of linear equations and
bilinear forms (expressions $\vec{x}\T\mat{A}\vec{y}$ where pairs of vectors interact
through a coefficient grid) worked heavily in fixed coordinate systems. They knew
equations could be redundant, but lacked any coordinate-free measure of a matrix's
true dimension, leaving the equality of row rank and column rank an unsolved puzzle.

At the time, determinants were the supreme algebraic tool; matrices were still viewed
merely as rectangular tables for calculating determinants. In 1879, Ferdinand Georg
Frobenius published his landmark paper \emph{Theorie der linearen Formen mit ganzen
Coefficienten}, formally defining the \textbf{rank} of a matrix as the size of its
largest non-vanishing \textbf{minor}. (A minor is simply the determinant of a square
submatrix formed by crossing out unwanted rows and columns.) Because the determinant
of any square matrix equals the determinant of its transpose, crossing out rows and
columns yields the exact same collection of non-zero minors whether you inspect
$\mat{A}$ or $\mat{A}\T$. Frobenius's determinantal definition settled algebraically
why row rank must strictly equal column rank.

Along a complementary track, the British mathematician James Joseph Sylvester
developed \textbf{invariant theory} (the study of algebraic properties that remain
unchanged under coordinate shifts), introducing the term ``nullity'' and publishing
his ``law of nullity'' in 1884:
\[
  \rank(\mat{A}\mat{B}) \;\ge\; \rank(\mat{A}) + \rank(\mat{B}) - n.
\]
Whereas Kronecker and Frobenius favored explicit determinantal criteria,
twentieth-century algebraists combined rank and nullity with linear maps to classify
operators into canonical forms (building on Camille Jordan's normal form).

Yet rank faces a modern computational challenge: exact algebraic rank is brittle under
floating-point roundoff. Consider the matrix $\begin{bmatrix} 1 & 1 \\ 1 & 1 + 10^{-16} \end{bmatrix}$.
Mathematically, its rows are independent and its rank is $2$. But in IEEE double-precision
floating-point arithmetic, $1 + 10^{-16}$ rounds to $1$, collapsing the computed rank to $1$.
Conversely, tiny roundoff noise can make a truly rank-$1$ physical matrix appear to have full rank!
Traditional elimination cannot distinguish between a tiny pivot that represents genuine physical
signal and one that represents roundoff garbage. This drove Gene Golub and William Kahan (1965)
to develop the singular value decomposition, replacing brittle integer rank with \emph{numerical rank}.
\end{historybox}

\begin{example}\label{ch16:rankfac-num}
Factor the rank-$1$ matrix from the opening:
\[
  \mat{A} = \begin{bmatrix} 1 & 2 \\ 2 & 4 \end{bmatrix}
  = \underbrace{\colvec{1 \\ 2}}_{\mat{C}}\;
    \underbrace{\begin{bmatrix} 1 & 2 \end{bmatrix}}_{\mat{F}}.
\]
The waist has width $1$. The single column $\colvec{1\\2}$ of $\mat{C}$ spans the
column space; the single row $\begin{bmatrix}1 & 2\end{bmatrix}$ of $\mat{F}$ spans
the row space. One number, $r=1$, and it is visibly the same on both sides.
\end{example}

\begin{connectionbox}
This equality parallels symmetry theorems elsewhere in mathematics, such as
Clairaut's theorem on the equality of mixed partial derivatives ($f_{xy} = f_{yx}$)
for smooth functions. In both settings, an underlying algebraic structure ensures
that evaluating operations across rows versus columns yields identical results.
\end{connectionbox}

\section{The four subspaces, paired by perpendicularity}

The rank $r$ determines the dimensions of all four fundamental subspaces.
Let $\mat{A}$ be $m\times n$ with rank $r$. Two subspaces live in the input space
$\R^n$ and two live in the output space $\R^m$:
\[
  \begin{array}{llcl}
  \text{in } \R^n: & \row(\mat{A}) & \text{dimension} & r,\\[2pt]
                   & \nullsp(\mat{A}) & \text{dimension} & n-r,\\[6pt]
  \text{in } \R^m: & \col(\mat{A}) & \text{dimension} & r,\\[2pt]
                   & \nullsp(\mat{A}\T) & \text{dimension} & m-r.
  \end{array}
\]
The dimensions already obey rank--nullity twice over: in $\R^n$ we have
$r + (n-r) = n$, and in $\R^m$ we have $r + (m-r) = m$. But the sharper fact is
\emph{geometric}: inside each space, the two subspaces are not just
complementary in dimension -- they are \textbf{orthogonal}. Every vector in one is
perpendicular to every vector in the other.

\begin{theorem}[Orthogonality of the fundamental subspaces]\label{ch16:orthog}
For any real $m\times n$ matrix $\mat{A}$,
\[
  \row(\mat{A}) \perp \nullsp(\mat{A}) \quad\text{in } \R^n,
  \qquad
  \col(\mat{A}) \perp \nullsp(\mat{A}\T) \quad\text{in } \R^m.
\]
What is more, each pair fills its whole space: $\row(\mat{A})$ and $\nullsp(\mat{A})$
together span $\R^n$, and $\col(\mat{A})$ and $\nullsp(\mat{A}\T)$ together span
$\R^m$. Each subspace is the \emph{orthogonal complement} of its partner. (An
\textbf{orthogonal complement} of a subspace $V$, written $V^\perp$, consists of
all vectors perpendicular to every vector in $V$.)
\end{theorem}

\begin{proof}
The first orthogonality is almost a restatement of what $\mat{A}\vec{x}=\vec0$
\emph{means}. Write $\mat{A}$ in terms of its rows
$\vec{r}_1\T,\dots,\vec{r}_m\T$. The equation $\mat{A}\vec{x}=\vec0$ says
precisely that every row dotted with $\vec{x}$ is zero:
\[
  \mat{A}\vec{x} = \vec0
  \quad\Longleftrightarrow\quad
  \inner{\vec{r}_i}{\vec{x}} = 0 \text{ for every } i.
\]
So a nullspace vector $\vec{x}$ is perpendicular to every row, and therefore to
every \emph{combination} of rows -- that is, to all of $\row(\mat{A})$. Hence
$\nullsp(\mat{A}) \perp \row(\mat{A})$.

For the dimensions: the two subspaces have dimensions $r$ and $n-r$, summing to
$n$, and orthogonal subspaces meet only at $\vec0$ (a vector perpendicular to
itself has length zero). Complementary dimension plus trivial intersection means
they together span $\R^n$, and being orthogonal, each is the orthogonal
complement of the other.

The second pair is the first pair applied to $\mat{A}\T$: the column space of
$\mat{A}$ is the row space of $\mat{A}\T$, and the left nullspace of $\mat{A}$ is
the nullspace of $\mat{A}\T$. Substituting $\mat{A}\T$ for $\mat{A}$ in what we
just proved gives $\col(\mat{A}) \perp \nullsp(\mat{A}\T)$ filling $\R^m$.
\end{proof}

\begin{intuitionbox}
Why must a nullspace vector be perpendicular to every row? Because
$\mat{A}\vec{x}$ is built out of dot products of $\vec{x}$ with the rows. For
$\mat{A}\vec{x}$ to be the zero vector, $\vec{x}$ has to give zero against
\emph{each} row -- it must point in a direction the rows cannot detect at all.
``Invisible to every row'' is exactly ``perpendicular to the row space.'' The
nullspace is the blind spot of the rows.
\end{intuitionbox}

\begin{figure}[H]
\centering
\begin{tikzpicture}[scale=1.0,>=Stealth]
  % input space R^n: row space (line) perp nullspace (line)
  \begin{scope}[shift={(-3.4,0)}]
    \draw[->,gray!60] (-2,0)--(2,0);
    \draw[->,gray!60] (0,-2)--(0,2);
    \draw[sourceblue,very thick] (-1.7,-1.7)--(1.7,1.7);
    \draw[targetgreen,very thick] (-1.7,1.7)--(1.7,-1.7);
    \draw[sourceblue,->,very thick] (0,0)--(1.1,1.1);
    \draw[targetgreen,->,very thick] (0,0)--(1.1,-1.1);
    % right angle mark
    \draw (0.22,0.22)--(0.44,0)--(0.22,-0.22);
    \node[sourceblue, above left] at (1.5,1.5) {\small $\row(\mat{A})$};
    \node[sourceblue, font=\scriptsize, below left] at (1.5,1.5) {$\dim = r$};
    \node[targetgreen] at (1.8,-1.15) {\small $\nullsp(\mat{A})$};
    \node[targetgreen,font=\scriptsize] at (1.8,-1.45) {$\dim = n-r$};
    \node at (0,-2.45) {\small input space $\R^n$};
  \end{scope}

  % mapping arrows between rooms
  \draw[->,very thick,sourceblue] (-1.9,1.5) to[bend left=25]
    node[midway,above,font=\small] {$\mat{A}$ (invertible, $r\to r$)} (1.9,1.5);
  \draw[->,thick,targetgreen,dashed] (-1.9,-1.2) to[bend right=20]
    node[midway,below,font=\small] {$\mat{A}\vec{x}_{\text{null}}=\vec{0}$} (3.2,-0.2);

  % output space R^m: col space perp left nullspace
  \begin{scope}[shift={(3.4,0)}]
    \draw[->,gray!60] (-2,0)--(2,0);
    \draw[->,gray!60] (0,-2)--(0,2);
    \draw[structurepurple,very thick] (-1.9,-0.8)--(1.9,0.8);
    \draw[quotientorange,very thick] (-0.8,1.9)--(0.8,-1.9);
    \draw[structurepurple,->,very thick] (0,0)--(1.4,0.59);
    \draw[quotientorange,->,very thick] (0,0)--(0.59,-1.4);
    % right angle mark in output space
    \draw (0.23,0.097)--(0.327,-0.133)--(0.097,-0.23);
    \node[structurepurple] at (1.9,1.15) {\small $\col(\mat{A})$};
    \node[structurepurple,font=\scriptsize] at (1.9,0.85) {$\dim = r$};
    \node[quotientorange] at (1.5,-1.35) {\small $\nullsp(\mat{A}\T)$};
    \node[quotientorange,font=\scriptsize] at (1.5,-1.65) {$\dim = m-r$};
    \fill[black] (0,0) circle (1.8pt);
    \node at (0,-2.45) {\small output space $\R^m$};
  \end{scope}
\end{tikzpicture}
\caption{The big picture of the Fundamental Theorem of Linear Algebra. In the input space
$\R^n$, the row space and nullspace are orthogonal complements ($r + (n-r) = n$). In the
output space $\R^m$, the column space and left nullspace are orthogonal complements
($r + (m-r) = m$). The matrix $\protect\mat{A}$ maps the row space \emph{one-to-one onto}
the column space, and annihilates the nullspace onto $\vec{0}$.}
\label{ch16:bigpicture}
\end{figure}

There is one more piece that turns these four subspaces into a single moving
picture, and it is the reason this whole circle of ideas earns the grand name
\emph{the Fundamental Theorem of Linear Algebra}.

\begin{theorem}[The action on the row space]\label{ch16:rowspace-iso}
The matrix $\mat{A}$ maps the row space \emph{one-to-one onto} the column space.
Restricted to the row space, $\mat{A}$ is invertible: distinct row-space vectors
go to distinct column-space vectors, and every column-space vector is hit exactly
once.
\end{theorem}

\begin{proof}
Both spaces have the same dimension $r$, so we only need that $\mat{A}$ is
one-to-one on the row space. Suppose two row-space vectors $\vec{u},\vec{v}$
satisfy $\mat{A}\vec{u}=\mat{A}\vec{v}$. Then $\mat{A}(\vec{u}-\vec{v})=\vec0$, so
$\vec{u}-\vec{v}$ lies in the nullspace. But $\vec{u}-\vec{v}$ is also a
difference of row-space vectors, hence in the row space. A vector in both the row
space and the nullspace is perpendicular to itself (Theorem
\ref{ch16:orthog}), so it is $\vec0$. Thus $\vec{u}=\vec{v}$: the map is
one-to-one on the row space. A one-to-one linear map between spaces of equal
dimension is onto as well.
\end{proof}

\begin{labbox}
\textbf{Mental check: decomposing an input vector.}
Because $\R^n = \row(\mat{A}) \oplus \nullsp(\mat{A})$, every input vector $\vec{x}\in\R^n$
splits uniquely into $\vec{x} = \vec{x}_{\text{row}} + \vec{x}_{\text{null}}$.
What is $\mat{A}\vec{x}$?
(Answer: $\mat{A}\vec{x} = \mat{A}\vec{x}_{\text{row}} + \mat{A}\vec{x}_{\text{null}}
= \mat{A}\vec{x}_{\text{row}} + \vec{0} = \mat{A}\vec{x}_{\text{row}}$. Because
$\mat{A}\vec{x}_{\text{null}} = \vec{0}$, all of the output comes from
$\vec{x}_{\text{row}}$, which $\mat{A}$ carries invertibly onto the column space.)
\end{labbox}

\begin{focusbox}
\textbf{The Fundamental Theorem of Linear Algebra, summarized.} Every real
$m\times n$ matrix splits both its spaces into orthogonal subspaces and acts
invertibly between them:
\begin{itemize}
  \item $\R^n = \row(\mat{A}) \oplus \nullsp(\mat{A})$, an orthogonal split into
        dimensions $r$ and $n-r$.
  \item $\R^m = \col(\mat{A}) \oplus \nullsp(\mat{A}\T)$, an orthogonal split into
        dimensions $r$ and $m-r$.
  \item $\mat{A}$ sends the row space \emph{invertibly onto} the column space, and
        maps the nullspace to $\vec0$.
\end{itemize}
The single number $r=\rank(\mat{A})$ controls every dimension in the diagram. This
is the structure the singular value decomposition will later make completely
explicit -- with orthonormal bases for all four subspaces at once.
\end{focusbox}

\begin{example}\label{ch16:four-numeric}
Let
\[
  \mat{A} = \begin{bmatrix} 1 & 2 & 3 \\ 2 & 4 & 6 \end{bmatrix}.
\]
Row $2$ is twice row $1$, so $r=1$ ($m=2$, $n=3$). The dimensions must be:
\[
  \dimn\row = 1,\quad \dimn\nullsp = n-r = 2,\quad
  \dimn\col = 1,\quad \dimn\nullsp(\mat{A}\T) = m-r = 1.
\]
\emph{Row space} ($\subset \R^3$): spanned by $\vec{r}=\colvec{1\\2\\3}$.
\emph{Nullspace} ($\subset \R^3$): solve $x_1 + 2x_2 + 3x_3 = 0$; the two special
solutions are $\colvec{-2\\1\\0}$ and $\colvec{-3\\0\\1}$. Check perpendicularity
to $\vec{r}$:
\[
  \inner{\vec r}{\colvec{-2\\1\\0}} = -2+2+0 = 0,\qquad
  \inner{\vec r}{\colvec{-3\\0\\1}} = -3+0+3 = 0. \checkmark
\]
\emph{Column space} ($\subset \R^2$): spanned by $\colvec{1\\2}$.
\emph{Left nullspace} ($\subset \R^2$): solve $\mat{A}\T\vec{y}=\vec0$, i.e.
$y_1 + 2y_2 = 0$ from every column; one solution direction is $\colvec{2\\-1}$,
and indeed $\inner{\colvec{1\\2}}{\colvec{2\\-1}} = 2-2 = 0$. \checkmark Both
perpendicular splits check out, and every dimension is governed by the lone
$r=1$.
\end{example}

\begin{trapbox}
The nullspace and the column space are \emph{not} a perpendicular pair, even
though both appear in the four-subspace list. The nullspace lives in $\R^n$
(input side); the column space lives in $\R^m$ (output side). They are in
different rooms and need not even have the same length vectors. The correct
partners are $\row \perp \nullsp$ (both in $\R^n$) and
$\col \perp \nullsp(\mat{A}\T)$ (both in $\R^m$). Pair by \emph{room}, not by name.
\end{trapbox}

\begin{walkbox}
A square $3\times 3$ matrix has rank $2$. Picture its four subspaces. The
nullspace is a line in $\R^3$; the row space is the plane perpendicular to that
line. On the output side, the column space is a plane in $\R^3$ and the left
nullspace is the line perpendicular to it. Now ask: is the input nullspace line
the same as the output left-nullspace line? Walk until you can produce a $3\times3$
matrix where they differ -- and one (symmetric) where they coincide.
\end{walkbox}

\section{Rank-one matrices: elementary outer products}

We have seen that rank measures the width of a matrix's waist. The narrowest
possible nonzero waist has width $1$. A matrix of rank $1$ is the simplest
nontrivial matrix there is -- and \emph{every} matrix is a sum of
rank-one matrices. They serve as elementary building blocks for general matrices
and play a central role in the singular value decomposition.

\begin{definition}\label{ch16:rank-one-def}
A nonzero matrix $\mat{M}$ has \textbf{rank one} if it can be written as an
\textbf{outer product}
\[
  \mat{M} = \vec{u}\,\vec{v}\T,
\]
where $\vec{u}\in\R^m$ and $\vec{v}\in\R^n$ are nonzero column vectors. Its
$(i,j)$ entry is $u_i v_j$.
\end{definition}

\begin{example}\label{ch16:outer}
With $\vec{u}=\colvec{1\\2\\3}$ and $\vec{v}=\colvec{4\\5}$,
\[
  \vec{u}\,\vec{v}\T
  = \colvec{1\\2\\3}\begin{bmatrix}4 & 5\end{bmatrix}
  = \begin{bmatrix} 4 & 5 \\ 8 & 10 \\ 12 & 15 \end{bmatrix}.
\]
Every row is a multiple of $\begin{bmatrix}4 & 5\end{bmatrix}$ (the multipliers
are $1,2,3$); every column is a multiple of $\colvec{1\\2\\3}$ (the multipliers
are $4,5$). Row space is a line, column space is a line: rank exactly $1$.
\end{example}

\begin{intuitionbox}
A rank-one matrix $\vec{u}\vec{v}\T$ acts on an input $\vec{x}$ via
\[
  (\vec{u}\vec{v}\T)\vec{x} = \vec{u}\,(\vec{v}\T\vec{x})
  = (\inner{\vec v}{\vec x})\,\vec{u}.
\]
The scalar product $\inner{\vec v}{\vec x}$ determines the magnitude, and the
output always points along the line spanned by $\vec{u}$. The vector $\vec{v}$
defines the linear functional that evaluates the input, while $\vec{u}$ specifies
the direction of the output in $\R^m$. Figure~\ref{ch16:rankone} illustrates
this two-step geometric evaluation.
\end{intuitionbox}

\begin{figure}[H]
\centering
\begin{tikzpicture}[scale=1.0,>=Stealth]
  % input side: x and the measuring stick v
  \begin{scope}[shift={(-3.5,0)}]
    \draw[->,gray!60] (-0.4,0)--(2.2,0);
    \draw[->,gray!60] (0,-0.4)--(0,2.2);
    \draw[->,targetgreen,very thick] (0,0)--(1.7,0.6) node[right]{$\vec{v}$};
    \draw[->,vecblue,very thick] (0,0)--(0.6,1.7) node[above]{$\vec{x}$};
    % projection drop of x onto v direction
    \draw[gray,dashed] (0.6,1.7)--(0.945,0.333);
    \fill[vecblue] (0.945,0.333) circle (1.1pt);
    \node[gray] at (1.55,-0.35) {\small overlap $\inner{\vec v}{\vec x}$};
    \node at (1.0,-0.9) {\small input space $\R^n$};
  \end{scope}
  % arrow
  \draw[->,thick] (-0.9,0.6)--(0.3,0.6) node[midway,above]{$\vec u\vec v\T$};
  % output side: everything lands on the ray through u
  \begin{scope}[shift={(2.0,0)}]
    \draw[->,gray!60] (-0.4,0)--(2.2,0);
    \draw[->,gray!60] (0,-0.4)--(0,2.2);
    \draw[structurepurple,very thick,dashed] (0,0)--(2.0,1.6);
    \draw[->,vecred,very thick] (0,0)--(1.25,1.0) node[below right]{$(\inner{\vec v}{\vec x})\,\vec{u}$};
    \node[structurepurple] at (2.05,1.95) {\small $\spn\{\vec u\}$};
    \node at (1.0,-0.9) {\small output: one ray};
  \end{scope}
\end{tikzpicture}
\caption{A rank-one matrix $\protect\vec{u}\protect\vec{v}\protect\T$ evaluated in two
stages. First, the dot product $\protect\inner{\vec v}{\vec x}$ produces a scalar.
Second, scaling $\protect\vec{u}$ by this scalar yields the output. Every output
lies in the one-dimensional subspace $\protect\spn\{\protect\vec{u}\}$.}
\label{ch16:rankone}
\end{figure}

\subsection{Every matrix is a sum of rank-one pieces}

Here is the bridge. Multiply $\mat{A}$ by a vector and you compute a linear
combination of $\mat{A}$'s columns. But there is a second, equally true way to
read a matrix product that reveals the fundamental rank-one building blocks.

Recall the rank factorization $\mat{A}=\mat{C}\mat{F}$ with waist $r$. A product
of an $m\times r$ matrix and an $r\times n$ matrix can be expanded as a sum over
the shared index: column $k$ of $\mat{C}$ times row $k$ of $\mat{F}$, summed.
Writing $\vec{c}_k$ for the columns of $\mat{C}$ and $\vec{f}_k\T$ for the rows
of $\mat{F}$,
\[
  \mat{A} = \mat{C}\mat{F}
  = \sum_{k=1}^{r} \vec{c}_k\,\vec{f}_k\T.
\]
Each term $\vec{c}_k\vec{f}_k\T$ is a rank-one matrix. So:

\begin{theorem}[Rank as a sum of rank-one pieces]\label{ch16:sum-rank-one}
A matrix has rank $r$ if and only if it can be written as a sum of $r$ rank-one
matrices, and $r$ is the \emph{smallest} number of rank-one matrices whose sum is
$\mat{A}$.
\end{theorem}

\begin{proof}[Why $r$ is the minimum]
If $\mat{A}=\sum_{k=1}^{p}\vec{u}_k\vec{v}_k\T$, then every column of $\mat{A}$ is
a combination of the $p$ vectors $\vec{u}_1,\dots,\vec{u}_p$ (each term
contributes a multiple of its $\vec{u}_k$). So the column space is spanned by
those $p$ vectors, giving $\rank(\mat{A})\le p$. Hence no sum of fewer than
$r=\rank(\mat{A})$ rank-one matrices can equal $\mat{A}$. The factorization above
achieves exactly $r$, so $r$ is the minimum.
\end{proof}

\begin{example}\label{ch16:sum-numeric}
Decompose a rank-$2$ matrix into two rank-one atoms. Take
\[
  \mat{A} = \begin{bmatrix} 1 & 0 & 2 \\ 0 & 1 & 3 \\ 1 & 1 & 5 \end{bmatrix}.
\]
Notice row$_3 = $ row$_1 + $row$_2$, so $r=2$. Columns $1$ and $2$,
$\vec{c}_1=\colvec{1\\0\\1}$ and $\vec{c}_2=\colvec{0\\1\\1}$, are a basis for the
column space. Column $3 = 2\vec{c}_1 + 3\vec{c}_2$. So the coefficient matrix
$\mat{F}$ (each column gives the recipe for the matching column of $\mat{A}$) is
\[
  \mat{F} = \begin{bmatrix} 1 & 0 & 2 \\ 0 & 1 & 3 \end{bmatrix},
  \qquad
  \mat{C} = \begin{bmatrix} 1 & 0 \\ 0 & 1 \\ 1 & 1 \end{bmatrix}.
\]
Now read off the two rank-one pieces -- column $k$ of $\mat C$ times row $k$ of
$\mat F$:
\[
  \mat{A}
  = \underbrace{\colvec{1\\0\\1}\begin{bmatrix}1 & 0 & 2\end{bmatrix}}_{\text{piece }1}
  + \underbrace{\colvec{0\\1\\1}\begin{bmatrix}0 & 1 & 3\end{bmatrix}}_{\text{piece }2}
  = \begin{bmatrix} 1 & 0 & 2 \\ 0 & 0 & 0 \\ 1 & 0 & 2 \end{bmatrix}
  + \begin{bmatrix} 0 & 0 & 0 \\ 0 & 1 & 3 \\ 0 & 1 & 3 \end{bmatrix}.
\]
Add them and you recover $\mat{A}$ exactly. Two atoms, because the rank is two.
\end{example}

\begin{figure}[H]
\centering
\begin{tikzpicture}[scale=0.9]
  % A as a solid block
  \draw[fill=sourceblue!12,draw=sourceblue,thick] (0,0) rectangle (1.8,1.8);
  \node at (0.9,0.9) {$\mat{A}$};
  \node[below=2pt] at (0.9,0) {\small rank $r$};
  \node at (2.4,0.9) {$=$};
  % layered rank-one atoms, drawn as a stack of thin slabs
  \begin{scope}[shift={(3.0,0)}]
    \foreach \i/\op in {0/30, 1/20, 2/12} {
      \begin{scope}[shift={(0.35*\i,0.35*\i)}]
        \draw[fill=quotientorange!\op,draw=quotientorange,thick] (0,0) rectangle (1.8,1.8);
      \end{scope}
    }
    \node at (1.25,1.25) {\small $\vec{c}_1\vec{f}_1\T$};
    \node at (0.55,0.55) {\small $+\,\cdots$};
  \end{scope}
  \node[below=2pt] at (4.6,0) {\small $r$ rank-one atoms};
  \node[structurepurple,above] at (4.6,2.6) {\small sum of $r$ outer products};
\end{tikzpicture}
\caption{Every matrix of rank $r$ can be written as a sum of $r$ rank-one matrices
$\protect\vec{c}_k\protect\vec{f}_k\T$ (Theorem \protect\ref{ch16:sum-rank-one}).
When the pieces are built from orthonormal vectors and ordered by weight
$\protect\sigma_1\ge\protect\sigma_2\ge\cdots>0$ (the singular value decomposition
in Chapter~29), dropping the smallest terms yields the optimal low-rank
approximation -- the mathematical engine of modern data compression.}
\label{ch16:atomstack}
\end{figure}

\begin{labbox}
By hand, write $\mat{A}=\begin{bmatrix}2 & 6 \\ 1 & 3\end{bmatrix}$ as a single
outer product $\vec{u}\vec{v}\T$. (Hint: every column is a multiple of
$\colvec{2\\1}$, so try $\vec{u}=\colvec{2\\1}$ and find the row $\vec{v}\T$ that
makes the columns $\colvec{2\\1}$ and $\colvec{6\\3}$ come out right. Answer:
$\vec{v}\T = \begin{bmatrix}1 & 3\end{bmatrix}$.) Then confirm $\rank(\mat A)=1$
by checking the second column is $3\times$ the first.
\end{labbox}

\begin{applicationbox}
\textbf{Low-rank approximation: how data gets compressed.} A grayscale image is a
matrix of pixel brightnesses -- maybe $1000\times 1000$, a million numbers. Most
real images are \emph{nearly} low rank: their information concentrates in a few
dominant directions. The singular value decomposition (covered in a later
chapter) writes any matrix as a sum of rank-one outer products ordered by importance
(the strength-ordered version of Figure \ref{ch16:atomstack}),
\[
  \mat{A} = \sigma_1 \vec{u}_1\vec{v}_1\T + \sigma_2 \vec{u}_2\vec{v}_2\T
          + \cdots + \sigma_r \vec{u}_r\vec{v}_r\T,
  \qquad \sigma_1 \ge \sigma_2 \ge \cdots > 0,
\]
and keeping only the first few terms gives the \emph{best} approximation of that
rank. Keep $50$ terms of a $1000\times 1000$ image and you store
$50\,(1000+1000)=100{,}000$ numbers instead of a million -- a $10\times$
compression -- while the picture stays recognizable. The same technique powers
recommender systems (a user-by-movie ratings matrix is low rank because tastes
cluster), latent-semantic analysis of text, and the parameter-efficient
fine-tuning of large neural networks (LoRA adds a trained low-rank update
$\mat{B}\mat{A}$ to a frozen weight matrix). Rank-one decompositions provide
the foundation for low-rank approximations across data science and numerical
computing.
\end{applicationbox}

\begin{recoverybox}
\textbf{Recap, and what each idea costs you.}
\begin{itemize}
  \item \textbf{Rank $=$ true reach.} The dimension the outputs fill -- equally,
        independent columns, independent rows, or pivots. \emph{Cost:} one
        elimination, $O(mn\cdot\min\{m,n\})$ arithmetic; then just count pivots.
  \item \textbf{Rank--nullity} (Theorem \ref{ch16:rank-nullity}):
        $\rank + \dimn\nullsp = n$. \emph{Cost:} free -- once you know the rank and
        the number of columns, the nullity needs no further work.
  \item \textbf{Row rank $=$ column rank} (Theorem \ref{ch16:row-col}): both count
        the width $r$ of the waist in $\mat{A}=\mat{C}\mat{F}$. \emph{Cost:} the
        same single elimination tells you both at once.
  \item \textbf{Four subspaces, two perpendicular pairs} (Theorem
        \ref{ch16:orthog}): $\row\perp\nullsp$ in $\R^n$,
        $\col\perp\nullsp(\mat{A}\T)$ in $\R^m$, all dimensions set by $r$.
        \emph{Cost:} the pivot columns and special solutions you already have hand
        you bases for all four.
  \item \textbf{Rank-one components} (Theorem \ref{ch16:sum-rank-one}): $\mat{A}$ is a
        sum of exactly $r$ outer products $\vec{c}_k\vec{f}_k\T$. \emph{Cost:} read
        them straight off the rank factorization; storing them costs $r(m+n)$
        numbers instead of $mn$ -- the whole point of low-rank compression.
\end{itemize}
One integer, $r$, paid for once, governs everything above.
\end{recoverybox}

\begin{connectionbox}
The rank--nullity conservation law $\rank(\mat{A}) + \dim\nullsp(\mat{A}) = n$ and
the outer-product expansion $\mat{A} = \sum_{k=1}^r \vec{c}_k \vec{f}_k\T$ give us a
complete structural inventory of what a matrix preserves and destroys. We can compress
low-rank matrices into $r(m+n)$ numbers instead of $mn$, forming the foundation of modern
data compression and machine learning.

Even so, two practical challenges remain. First, pure algebraic rank is brittle under
floating-point roundoff: a matrix can be mathematically full rank while sitting within
$10^{-16}$ of a lower-rank matrix. (In Chapter~29, the Singular Value Decomposition
will address this by replacing the integer rank with a continuous spectrum
of singular values.)

Second, while Theorem~\ref{ch16:orthog} establishes that the row and column spaces are
perpendicular to their partner nullspaces, the basis vectors \emph{within} each subspace
may be non-orthogonal. In Chapter~17, we launch Part~IV by developing orthogonal
and orthonormal coordinate systems, where coordinates decouple and projections
onto subspaces simplify significantly.
\end{connectionbox}

\section{Exercises}

\begin{exercise}\label{ch16:ex-warmup}
Find the rank of each matrix by reducing to echelon form and counting pivots.
\[
  \text{(a) } \begin{bmatrix} 3 & 6 \\ 1 & 2 \end{bmatrix}, \quad
  \text{(b) } \begin{bmatrix} 1 & 2 & 0 \\ 0 & 0 & 1 \\ 1 & 2 & 1 \end{bmatrix}, \quad
  \text{(c) } \begin{bmatrix} 1 & 1 & 1 \\ 1 & 2 & 3 \\ 1 & 3 & 5 \end{bmatrix}.
\]
For each, state the nullity using rank--nullity, then verify by finding the
special solutions.
\end{exercise}

\begin{exercise}\label{ch16:ex-rn-predict}
A matrix $\mat{A}$ is $4\times 7$ and has rank $3$. Without seeing the entries,
state: the dimension of the nullspace, the dimension of the column space, the
dimension of the left nullspace, and the dimension of the row space. Which two of
these subspaces live in $\R^4$, and which two in $\R^7$?
\end{exercise}

\begin{exercise}\label{ch16:ex-perp}
For $\mat{A}=\begin{bmatrix} 1 & 2 & 2 \\ 2 & 4 & 5 \end{bmatrix}$, find a basis
for the row space and a basis for the nullspace. Verify by direct dot product
that every nullspace basis vector is perpendicular to every row space basis
vector. (Determine the dimension of the nullspace from the rank first.)
\end{exercise}

\begin{exercise}\label{ch16:ex-outer}
Write $\mat{A}=\begin{bmatrix} 6 & 4 & 2 \\ 9 & 6 & 3 \end{bmatrix}$ as a single
outer product $\vec{u}\vec{v}\T$. Then write its rank. Could you have known the
rank just by glancing at the rows? Explain in one sentence.
\end{exercise}

\begin{exercise}\label{ch16:ex-sum}
Decompose $\mat{A}=\begin{bmatrix} 1 & 1 & 2 \\ 1 & 2 & 3 \\ 0 & 1 & 1 \end{bmatrix}$
into a sum of rank-one matrices using a rank factorization $\mat{A}=\mat{C}\mat{F}$.
First find the rank (look for a column dependency), then build $\mat{C}$ from a
basis of the column space and $\mat{F}$ from the recipes for each column, and
finally write $\mat{A}=\sum_k \vec{c}_k\vec{f}_k\T$. Check by adding the pieces.
\end{exercise}

\begin{exercise}\label{ch16:ex-rankone-property}
Let $\mat{M}=\vec{u}\vec{v}\T$ be a nonzero rank-one $n\times n$ matrix. Show that
$\mat{M}^2 = (\inner{\vec v}{\vec u})\,\mat{M}$. Conclude that if
$\inner{\vec v}{\vec u}=0$ then $\mat{M}^2 = \mat{0}$, even though
$\mat{M}\ne\mat{0}$. Give a concrete $2\times 2$ example with $\mat{M}^2=\mat0$.
\end{exercise}

\begin{exercise}\label{ch16:ex-rank-bound}
Prove that for any matrices $\mat{A}$ (size $m\times k$) and $\mat{B}$ (size
$k\times n$), $\rank(\mat{A}\mat{B}) \le \min\{\rank(\mat{A}),\rank(\mat{B})\}$.
(Hint: every column of $\mat{A}\mat{B}$ is a combination of the columns of
$\mat{A}$; then apply the same idea to rows, or to the transpose.) Where does this
say the ``waist'' of a product cannot be wider than either factor's waist?
\end{exercise}

\begin{exercise}\label{ch16:ex-rankone-sum-trap}
Find two rank-one matrices whose sum has rank $1$, and two whose sum has rank $2$.
Then find two nonzero rank-one matrices whose sum is the zero matrix (rank $0$).
What does this show about the claim ``rank of a sum equals sum of ranks''? (It is
false; show the correct statement is $\rank(\mat{A}+\mat{B})\le
\rank(\mat{A})+\rank(\mat{B})$ and explain the inequality in one line using column
spaces.)
\end{exercise}

\begin{exercise}\label{ch16:ex-rowcol-proof}
Using the rank factorization $\mat{A}=\mat{C}\mat{F}$ with waist $r=$ column rank,
explain in your own words why the row space of $\mat{A}$ is spanned by the rows of
$\mat{F}$, and hence why row rank $\le$ column rank. Then explain why applying the
same statement to $\mat{A}\T$ closes the loop to equality. (Aim for a paragraph a
classmate could follow without the figure.)
\end{exercise}

\begin{exercise}\label{ch16:ex-rowspace-iso}
Let $\mat{A}=\begin{bmatrix} 1 & 0 & 1 \\ 0 & 1 & 1 \end{bmatrix}$ (rank $2$,
$n=3$). Take the two row-space basis vectors $\vec{r}_1=\colvec{1\\0\\1}$ and
$\vec{r}_2=\colvec{0\\1\\1}$ and compute $\mat{A}\vec{r}_1$ and $\mat{A}\vec{r}_2$.
Verify the outputs are independent (so $\mat{A}$ is one-to-one on the row space,
as Theorem~\ref{ch16:rowspace-iso} promises). Then take the nullspace vector
$\vec{n}=\colvec{1\\1\\-1}$, confirm $\mat{A}\vec{n}=\vec0$, and note how the row
space gets mapped one-to-one while the nullspace is annihilated.
\end{exercise}

\begin{exercise}\label{ch16:ex-challenge}
(Challenge.) Suppose $\mat{A}$ is $n\times n$ with $\rank(\mat{A})=r<n$. Prove that
$\mat{A}$ has a nonzero nullspace, and use rank--nullity to say exactly its
dimension. Then prove the converse direction of the invertibility dictionary: an
$n\times n$ matrix is invertible if and only if its rank is $n$. (Connect ``rank
$n$'' to ``trivial nullspace'' to ``columns independent'' to ``invertible.'')
\end{exercise}
